\chapter{Deployment and Project Artifacts}

\section{Purpose}
This appendix summarizes the deployment-facing and documentation-facing artifacts used in the SARS project. The content is written directly in the report so that the thesis compiles reliably on Overleaf without requiring external artifact files.

\section{Representative Backend Dependency Set}
\begin{lstlisting}[style=sarscode,caption={Representative backend dependency declaration excerpt}]
fastapi
uvicorn
sqlalchemy
pydantic
pydantic-settings
python-jose
passlib
bcrypt
python-multipart
chromadb
google-generativeai
\end{lstlisting}

\section{Representative Frontend Package Set}
\begin{lstlisting}[style=sarscode,caption={Representative frontend package declaration excerpt}]
{
  "name": "sars-frontend",
  "private": true,
  "dependencies": {
    "axios": "...",
    "react": "...",
    "react-dom": "...",
    "react-router-dom": "..."
  }
}
\end{lstlisting}

\section{Project Artifact Summary}
\begin{table}[H]
\centering
\caption{Key Project Artifacts Preserved During Development}
\begin{tabularx}{\textwidth}{>{\raggedright\arraybackslash}p{0.28\textwidth} >{\raggedright\arraybackslash}X}
\toprule
\textbf{Artifact} & \textbf{Purpose} \\
\midrule
Goal manifest & Captured implementation scope, milestone readiness, and system completeness notes \\
Verification notes & Recorded milestone checks, bug fixes, and validation outcomes during development \\
Project report structure & Guided the chapter-wise layout and screenshot planning of the thesis \\
Overleaf report package & Preserved a portable version of the thesis source for cloud compilation \\
Screenshot placement map & Mapped raw UI screenshots to thesis-safe filenames and report sections \\
\bottomrule
\end{tabularx}
\end{table}

\section{Deployment Notes}
The implemented system is organized for a practical local deployment model:
\begin{itemize}[leftmargin=*]
\item the backend runs as a FastAPI application,
\item the frontend runs as a React client,
\item the database uses SQLite for simple project-scale deployment,
\item the vector store uses persistent local Chroma storage,
\item AI-backed functions depend on valid external API configuration.
\end{itemize}

\section{Environment Configuration Summary}
\begin{table}[H]
\centering
\caption{Representative Runtime Configuration Inputs}
\begin{tabularx}{\textwidth}{>{\raggedright\arraybackslash}p{0.24\textwidth} >{\raggedright\arraybackslash}X}
\toprule
\textbf{Configuration Item} & \textbf{Role in the System} \\
\midrule
\texttt{SECRET\_KEY} & Used for token signing and authentication integrity \\
\texttt{DATABASE\_URL} & Points the backend to the SQLite database file \\
\texttt{GEMINI\_API\_KEY} & Enables extraction, embeddings, and advisory generation \\
\texttt{CHROMA\_DB\_PATH} & Defines persistent vector-store location \\
\texttt{TEACHER\_INVITE\_CODE} & Protects teacher onboarding from uncontrolled access \\
\bottomrule
\end{tabularx}
\end{table}

\section{Interpretation}
These deployment and artifact summaries strengthen the engineering quality of the report. They show that the project was not only implemented functionally, but also packaged, verified, and documented in a way that supports demonstration, review, and future extension.
